\documentclass[11pt]{article}
\usepackage[margin=1in]{geometry}
\usepackage{listings}
\usepackage{xcolor}
\usepackage{enumitem}
\usepackage{amsmath}

\lstset{
  language=Java,
  basicstyle=\ttfamily\small,
  keywordstyle=\bfseries\color{blue!70!black},
  commentstyle=\color{green!50!black},
  stringstyle=\color{red!60!black},
  showstringspaces=false,
  frame=single,
  framesep=6pt,
  xleftmargin=6pt,
  xrightmargin=6pt,
  aboveskip=10pt,
  belowskip=10pt,
}

\pagestyle{empty}

\begin{document}

\begin{center}
  {\Large \textbf{CSCI-UA 102 --- Quiz 4}} \\[6pt]
  {\large Trees} \\[4pt]
  Name: \underline{\hspace{2.5in}} \quad NetID: \underline{\hspace{1.5in}}
\end{center}

The \textbf{height} of a node $p$ in a tree is the number of edges on the longest downward path from $p$ to a leaf (external node). A leaf has height 0. \noindent\textbf{Write an efficient, recursive method:}
\begin{lstlisting}
static <E> int sumHeight(Tree<E> tree)
\end{lstlisting}
that calculates the sum of the heights of every node in the tree. \noindent\textit{Hint: use an auxiliary variable.}
You may use the following interfaces and assume working implementations exist.
\begin{lstlisting}
public interface Position<E> {
    E getElement();
}
public interface Tree<E> {
    Position<E> root();
    Iterable<Position<E>> children(Position<E> p);
    boolean isExternal(Position<E> p);
    boolean isEmpty();
}
\end{lstlisting}


\newpage

{\color{blue}
\section*{Solution}

\begin{lstlisting}[basicstyle=\ttfamily\small\color{blue},
  keywordstyle=\bfseries\color{blue!70!black},
  commentstyle=\color{blue!50}]
static <E> int sumHeight(Tree<E> tree) {
    if (tree.isEmpty()) return 0;
    return subtreeHeightAndSum(tree, tree.root())[1];
}

// Returns {height, sumOfHeights} for the subtree at p.
static <E> int[] subtreeHeightAndSum(Tree<E> tree,
                                      Position<E> p) {
    int height = 0;
    int sum = 0;
    for (Position<E> child : tree.children(p)) {
        int[] childResult = subtreeHeightAndSum(tree, child);
        height = Math.max(height, 1 + childResult[0]);
        sum += childResult[1];
    }
    sum += height;                // add this node's height
    return new int[]{height, sum};
}
\end{lstlisting}

\textbf{Time complexity:} $O(n)$ --- each node is visited exactly once, with $O(c_p)$ work per node.

\medskip

\textbf{Why return two values?} We need the height (to recurse) \textit{and} the sum of heights (the answer). Returning \texttt{\{height, sumOfHeights\}} lets us compute both in a single bottom-up pass. Without this, computing the height for each node independently would be $O(n)$ per node, giving $O(n^2)$ total.
}

\end{document}
